\section{Stability-Based Selection}
\label{sec:selection}
Having established the normalized decision matrix $\mathbf{X}$ and the semantic definition of criteria, the core decision problem is now addressed by selecting the most robust maintenance schedule.

\subsection{Criteria Importance Weights}

The Weighted OWA (WOWA) operator requires a criteria importance vector $\mathbf{p} = (p_1, \dots, p_M)$ that dictates the ``potential influence'' of each criterion, independent of the realization of values. To examine how different organizational priorities affect the selection outcome, two contrasting importance configurations are considered (Table~\ref{tab:importance_vectors}): a \textit{Baseline} scenario reflecting standard operational priorities, and a \textit{Sustainability} scenario representing a strategic emphasis on environmental and social criteria.

\begin{table}[htbp]
\centering
\caption{Criteria importance vector configurations representing distinct operational priorities.}
\label{tab:importance_vectors}
\begin{adjustbox}{center}
\begin{tabular}{lccccccc}
\toprule
\textbf{Scenario} & \rotatebox{60}{\textbf{Rel.\ Stress}} & \rotatebox{60}{\textbf{Crit.\ Overlap}} & \rotatebox{60}{\textbf{Risk Volat.}} & \rotatebox{60}{\textbf{Rob.\ Slack}} & \rotatebox{60}{\textbf{Ski Season}} & \rotatebox{60}{\textbf{Eco-Stress}} & \rotatebox{60}{\textbf{Soc. Weekend}} \\
\midrule
Baseline & 0.10 & 0.25 & 0.10 & 0.20 & 0.10 & 0.10 & 0.15 \\
Sustainability & 0.20 & 0.05 & 0.10 & 0.15 & 0.10 & 0.25 & 0.15 \\
\bottomrule
\end{tabular}
\end{adjustbox}
\end{table}

The \textit{Baseline} scenario is designed to reflect standard transmission system operator priorities. A total importance mass of 0.65 is assigned to the four technical criteria, maintaining the priority of security of supply. Within this block, \textit{Max Critical Overlap} is assigned 0.25 and \textit{Robustness Slack} 0.20, because both metrics encode direct exposure to simultaneous critical interventions and the available operational buffer. \textit{Reliability Stress} and \textit{Risk Volatility} are each assigned 0.10, since both characterize technical performance at an aggregate level rather than through direct feasibility margins. The remaining 0.35 is assigned to the socio-environmental criteria, with \textit{Social Weekend Risk} receiving 0.15 to reflect the practical importance of public acceptability, while \textit{Ski Season Risk} and \textit{Eco-Stress} remain at 0.10 each.

The \textit{Sustainability} scenario is designed to represent a strategic setting where environmental compliance and social license are elevated to primary objectives without eliminating technical safeguards. The total importance mass is split evenly between technical criteria (0.50) and socio-environmental criteria (0.50). Within the socio-environmental block, \textit{Eco-Stress} is assigned 0.25 to reflect the possibility that environmental constraints become a binding driver of admissible schedules, while \textit{Social Weekend Risk} is assigned 0.15 to represent the relevance of public acceptance and local disruption. \textit{Ski Season Risk} is set to 0.10, capturing seasonal exposure as a material contextual factor. Within the technical block, \textit{Reliability Stress} is assigned 0.20 to maintain the priority of security, \textit{Robustness Slack} 0.15 to maintain reserve-based feasibility, \textit{Risk Volatility} 0.10 to retain sensitivity to operational instability, and \textit{Max Critical Overlap} 0.05 to keep simultaneous criticality under control without overriding the sustainability-driven objectives.

\subsection{Stability-Based WOWA Selection in a Unified Weight Space}

For a fixed importance vector and decision matrix $\mathbf{X}$, Theorem~\ref{thm:wowa} yields a unified weight-space representation of WOWA aggregations in which all alternatives are evaluated under a single common vector $\mathbf{W}$. Without this result, each alternative would be evaluated through its own cumulative-importance partition and its own induced weight vector, so equality between two WOWA scores would not define a hyperplane in a common simplex. Theorem~\ref{thm:wowa} removes this obstacle by refining all alternative-specific partitions into a shared grid, thereby restoring the linear geometry required for our exact stability analysis. In line with the classical literature on weight stability in additive MCDM methods \cite{mareschal1988weight}, the geometry of this common weight space is exploited to select the most robust solution, i.e., the alternative whose dominance is preserved under the widest range of weight variations.

The space of admissible weights is the unit simplex $\Delta^{K-1}$. For any pair of alternatives $a_i$ and $a_j$, the condition of indifference (same score) defines a hyperplane:
\[
H_{ij} = \{ \mathbf{W} \in \Delta^{K-1} \mid \mathbf{W}^T(\mathbf{z}_i - \mathbf{z}_j) = 0 \}
\]
The collection of $\binom{N}{2}$ indifference hyperplanes partition the simplex into convex polytopes, each corresponding to a specific ranking. For the choice problem (identifying the single best alternative) only $N-1$ hyperplanes are relevant, those comparing the top candidate against each other alternative.

As discussed earlier, the orness provides an interpretable summary of the decision-maker's attitude, but a single orness value does not uniquely determine the weight vector (see Figure~\ref{fig:orness_simplex}). Furthermore, a precise point estimate is often unrealistic. A more robust approach is to specify an \textit{orness interval} $[\alpha_{\min}, \alpha_{\max}]$, which defines a slice of the simplex $\Delta^{K-1}$. Within this feasible region, the goal is to find the most stable solution.

While the volume of a stability region is a common proxy for robustness, it can be misleading: a region may have large volume but be arbitrarily close to an indifference hyperplane. A more rigorous metric is the \textit{stability radius} $\rho$, defined as the radius of the largest Euclidean ball inscribed within the feasible region. This quantifies the minimum perturbation in weights required to alter the decision.

Therefore, to find the most stable choice, $N$ independent Linear Programs are solved (one per candidate $a_i$), maximizing the inscribed ball radius $\rho_i$ subject to dominance, simplex, and orness constraints. This geometric approach extends earlier LP-based sensitivity analyses for additive MCDM methods, where mathematical programming is used to determine the minimum modification of the weights required to make a given alternative ranked first \cite{wolters1995novel}. 

Since the weights lie on the hyperplane $\sum_{k=1}^K W_k = 1$ rather than in full $\mathbb{R}^K$, distances must be measured within this $(K-1)$-dimensional surface. The vector $\mathbf{1} = (1, \dots, 1)^T$ is normal to this hyperplane, so any admissible displacement of $\mathbf{W}$ must be orthogonal to $\mathbf{1}$. Standard Euclidean distance in $\mathbb{R}^K$ would underestimate the margin to constraint boundaries, as it measures along directions that leave the feasible surface. To compute distances correctly, vectors are projected via $\Pi \in \mathbb{R}^{K \times K}$. For any vector $\mathbf{v} \in \mathbb{R}^K$, the projection removes its component along $\mathbf{1}$:
\[
\Pi\mathbf{v} = \mathbf{v} - \frac{\mathbf{v}^T\mathbf{1}}{\|\mathbf{1}\|^2} \mathbf{1} = \mathbf{v} - \frac{1}{K}\left(\sum_{j=1}^K v_j\right) \mathbf{1}
\]

For any linear constraint $\mathbf{n}^T \mathbf{W} \ge b$, the distance from $\mathbf{W}$ to the boundary $\mathbf{n}^T \mathbf{W} = b$, measured along the hyperplane, is $(\mathbf{n}^T \mathbf{W} - b)/\|\Pi\mathbf{n}\|_2$. A ball of radius $\rho$ centered at $\mathbf{W}$ remains inside the halfspace if and only if $\mathbf{W}$ is at least $\rho$ away from the boundary, yielding the margin condition $\mathbf{n}^T \mathbf{W} \ge b + \rho \cdot \|\Pi\mathbf{n}\|_2$. 

Let $\boldsymbol{ r }=( r _1,\dots, r _K)$, with $ r _k = (K-k)/(K-1)$, denote the orness coefficients. Then, applying this to each constraint type:
\begin{itemize}
    \item \textbf{Simplex faces} ($W_k \ge 0$ for each $k$): Normal $\mathbf{e}_k$ \com{(the $k$-th canonical basis vector)}{esto no sé así está claro}, boundary $W_k = 0$. Yields \eqref{eq:lp_faces}.
    \item \textbf{Orness bounds}: Normal $\pm\boldsymbol{ r }$, boundaries $\boldsymbol{ r }^T \mathbf{W} = \alpha_{\min}$ and $\boldsymbol{ r }^T \mathbf{W} = \alpha_{\max}$. Yields \eqref{eq:lp_orness_margin}.
    \item \textbf{Dominance}: Normal $(\mathbf{z}_i - \mathbf{z}_j)$, boundary $\mathbf{W}^T(\mathbf{z}_i - \mathbf{z}_j) = 0$. Yields \eqref{eq:lp_dominance}.
\end{itemize}

Therefore, the Linear Problem for candidate $a_i$ is:
\begin{align}
    \textbf{Maximize } & \rho_i \label{eq:lp_objective}\\
    \textbf{Subject to: } & \notag \\
    & \sum_{k=1}^K W_k = 1, \quad W_k \ge 0 \quad \forall k \label{eq:lp_simplex}\\
    & \alpha_{\min} \le \boldsymbol{ r }^T \mathbf{W} \le \alpha_{\max} \label{eq:lp_orness}\\
    & W_k \ge \rho_i \cdot \|\Pi \mathbf{e}_k\|_2 \quad \forall k=1,\dots,K \label{eq:lp_faces}\\
    & \boldsymbol{ r }^T \mathbf{W} - \alpha_{\min} \ge \rho_i \cdot \|\Pi\boldsymbol{ r }\|_2, \quad \alpha_{\max} - \boldsymbol{ r }^T \mathbf{W} \ge \rho_i \cdot \|\Pi\boldsymbol{ r }\|_2 \label{eq:lp_orness_margin}\\
    & \mathbf{W}^T(\mathbf{z}_i - \mathbf{z}_j) \ge \rho_i \cdot \|\Pi (\mathbf{z}_i - \mathbf{z}_j)\|_2 \quad \forall j \neq i \label{eq:lp_dominance}
\end{align}

Constraints from \eqref{eq:lp_simplex} to \eqref{eq:lp_orness} define the feasible polytope, namely the simplex slice satisfying the orness bounds. Constraints from \eqref{eq:lp_faces} to \eqref{eq:lp_dominance} ensure that the inscribed ball of radius $\rho_i$ fits within this region.

If this linear problem is feasible, then alternative $a_i$ can be the optimal choice for some weight vector. The largest $\rho^*$ is selected as the most robust decision.
